\subsection{Neural dynamic capital policies and full unilateral deviations}\label{sec:dynamic-computation}
For two firms, we retain the capital game with $A=1$, $P_0=2$, $\delta=.1$, $\sigma_K=.15$, $\rho=.04$, $T=3$, investment in $[0,.8]$, and capital in $[.05,1.5]^2$. Terminal or joint-domain exit pays $.5K_i$. Baseline adjustment coefficients are $b_1=b_2=1$; the earlier asymmetric experiment has $b_2=1.2$. The principal neural grid has $17\times17$ states, 30 time steps, and 11 actions per player. Positive cubature and interpolation define the finite game. Market normalizations $M_2=1$ and $M_2=2$ are distinct economies, not evidence for an uncomputed large-population limit.

At each time and state, the R7 calculation trains separate critics and actors. It constructs both full $11\times11$ continuation-payoff matrices. The training labels are a pure one-step equilibrium when one exists and a minimum-maximum-defect pair otherwise. The actor's own objective is differentiated with the rival action and continuation critic fixed. Thus this study actually trains neural strategic policies, but it also pays for an exhaustive joint-action training oracle. Raw policies, three-by-three shortlist policies, and any corrected policies are distinct stored arrays.

After freezing the entire profile, an independent backward optimization permits every own action at every state and time against the fixed rival policy. Its value minus the profile payoff is the unilateral exploitability map. This is not a one-period test with the approximate continuation held fixed. The raw and final profiles meet the $.1$ absolute-profit target in every principal R7 case; the full-oracle guard is unused in those cases.
\begin{table}[htbp]\centering
\caption{R7 neural dynamic-game profiles and full best responses}
\begin{small}\begin{tabular}{lrrrrrr}
\toprule
Market & Seed & Raw max & Firm 1 & Firm 2 & Missing & Seconds \\
\midrule
1 & 11 & 0.0259 & 0.0164 & 0.0252 & 3 & 52.43 \\
2 & 11 & 0.0241 & 0.0255 & 0.0163 & 1 & 56.13 \\
1 & 29 & 0.0300 & 0.0288 & 0.0296 & 0 & 52.44 \\
2 & 29 & 0.0626 & 0.0624 & 0.0547 & 1 & 57.23 \\
1 & 47 & 0.0190 & 0.0081 & 0.0183 & 0 & 52.53 \\
2 & 47 & 0.0101 & 0.0074 & 0.0113 & 0 & 56.98 \\
\bottomrule\end{tabular}\end{small}
\par\smallskip\begin{minipage}{.96\linewidth}\footnotesize All bounds cover every stored finite state, time, and unilateral action, including deviations at the reported missing pure stages. The full-oracle guard is unused in these principal runs. Profit losses are absolute units; market normalizations are different economies.\end{minipage}
\end{table}

Five state-time nodes across the six neural runs lack a pure one-step equilibrium for the fitted continuation matrices. Their locations and defects are reported in the supplement. A minimum-defect selection at those nodes is not called an exact equilibrium. Nevertheless, the complete dynamic best-response bound remains valid for the resulting profile: its proof does not assume that every fitted one-step game has a pure solution. Moreover, for fixed Markov rivals, randomized or history-dependent own deviations cannot improve on the finite Markov best-response value, because each conditional action mixture is bounded by its largest pure-action payoff. The calculation therefore checks such unilateral deviations as well; it does not compute a mixed-strategy equilibrium profile or a continuous-action game.

R8 adds actor-free neural continuation fitting and classical backward selection for both market normalizations. The neural comparator allocates 480 critic Adam updates per time level rather than R7's 320 critic plus 160 actor updates, with the same L-BFGS cap. Counts and timings reveal the difference between nominal update matching and matched wall time. Both methods receive their own independent full dynamic best responses. The earlier six classical game cases, asymmetric costs, alternative tie rules, boundary frequencies, coarse counterexample, and finer-grid results remain in the supplement and raw record.
\begin{table}[htbp]\centering
\caption{R8 actor-free dynamic-game comparisons}
\begin{small}\begin{tabular}{lrrrrrr}
\toprule
Method & Market & Seed & Firm 1 & Firm 2 & Train (s) & Verify (s) \\
\midrule
Classical & 1 & 11 & 0.0000 & 0.0000 & 0.22 & 6.40 \\
Classical & 2 & 11 & 0.0000 & 0.0000 & 0.23 & 6.29 \\
Direct & 1 & 11 & 0.0223 & 0.0273 & 37.43 & 6.09 \\
Direct & 1 & 29 & 0.0199 & 0.0282 & 37.68 & 6.18 \\
Direct & 1 & 47 & 0.0223 & 0.0409 & 36.03 & 6.26 \\
Direct & 2 & 11 & 0.0091 & 0.0694 & 40.43 & 5.75 \\
Direct & 2 & 29 & 0.0306 & 0.0185 & 43.07 & 6.34 \\
Direct & 2 & 47 & 0.0206 & 0.0358 & 40.47 & 6.08 \\
\bottomrule\end{tabular}\end{small}
\par\smallskip\begin{minipage}{.96\linewidth}\footnotesize The direct method fits two continuation networks with the same class and 480 critic Adam updates, versus R7's 320 critic and 160 actor updates. Classical selection uses no neural approximation. Every method receives a new full dynamic best response. R7 and R8 timings were measured in their recorded environments; cross-environment ratios are not a hardware-controlled speed experiment.\end{minipage}
\end{table}

The finite-game result is a numerically evaluated approximate equilibrium with a specified profit-unit bound. It is not inferred from stationarity of a summed actor loss or the static $1/64$ diagnostic. State interpolation, continuous investment, diffusion timing, and continuously monitored exits remain separate terms in a continuous-game error account.
